\section{Selección de Herramienta para Métricas de Consumo}
\subsection{Muestra de Videos Para la fase de Experimentacion 1}
% Fase de la Experimentacion 2
% ENTORNO PRUEBAS EXP2:
%   + Caracteristicas del video
%       - Duracio 10s
%       - CRF 0
%       - Calidad 4k
%       - Bit da cada video [](TABLA)
%       - Codificador base x264(FFMPEG PROBE)(TABLA)
%   + Lugar de experimentacion --> Maquina de fase 1
% Seleccion de codificadorres: 
%  + Basandose en tabla \ref{tab:nvidia_matrix_capacity} son principalmente H.264 y H.265,

\subsection{Metricas de consumo: Transcodificacion usando 1 contendor sobre CPU+GPU}
%  +++ COMPARATIVA HERRAMIENTAS +++
% ++++++ EXPERIMENTO 1 +++++++++
% - Comparativa de EM - Kepler - Scaphandre || CPU + GPU
%   - Un aspecto negativo de Energy Meter es que no permite metricas aisladas por proceso, por tanto no es una opcion valida como herramienta final.
%   - Schapahndre permite obtener metricas de consumo detallada por proceso, en cuanto al uso de CPU, pero no es capaz de tomar metricas de la GPU. La toma de 
%     metricas de consumo energetico es fundamental para el proyecto, ya que la codificacion de video se lleva a cabo tambien en las GPU. Por este motivo
%     se descarta su eleccion como herramienta final
%   - Escogemos Kepler porque mide GPU ya que scp no mide gpu y em no mide por proesos ni permite segmentar procesos por cores de cpu

\subsection{Metricas de consumo: Transcodificacion usando 1 contendor sobre CPU}
% -Comparativa empleando codificacion unicamente sobre CPU // Kepler vs SCP vs EM
% -En la siguente fase se lanzara un proceso unicamente utilizando CPU para comprobar la precision de ambas herramientas, teniendo en cuenta que 
%   scp no es capaz de medir la GPU, como se ha visto en la prueba anterior.
% -Los resultados obtendios muestran XXXXX 

\subsection{Metricas de consumo: Transcodificacion usando 1 contenedor con 3 cores de CPU limitados}
% ++++++ EXPERIMENTO 2 +++++++++
% - Compartiva aislando Cores - Kepler vs SCP vs EM || CPU
%   
%   - En el siguiente experimento se medira la precision de ambas herramientas para medir procesos con limitacion de CPU , sin hacer uso 
%       de la acceleracion mediente GPU.
%     Para este caso particualar, se limitaran las CPUs a 3 al lanzar el contenedor de Docker empleando las flags CODE(--cpus="3" --cpuset-cpus=1-3).
%   - Para esta comparativa se plantearan dos escenarios, uno en el que o se limita el numero de CPUs y otro en el que se limitaran las CPUs a 3. Para cada
%       escenario se lanzaran 10 iteraciones de codificacion y toma de metricas y posteriormete se compara entre ellas con el fin de ver si exite discrepancia
%    entre la cantidad de energia que la herramienta mide del proceso con y sin limitaciones respecto a la energia total del sistema. 
%   -  Se observa XXXX

\subsection{Metricas de consumo: Transcodificacion usando 2 contenedores simultaneos con limite de 3 cores cada uno}
% ++++++ EXPERIMENTO 3 +++++++++
% - Contenedores Limitados a 3 CORES // EM - Kepler - Scaphandre || CPU 
%   - Para finalizar la comparativa de herramientas, se lleva a acabo un ultimo experimento en el que cual Kepler y SCP tendran que tomar metricas de dos contenedores
%   los cuales emplean codificacion sobre CPU con limitaciones a 3 en el numero de cores. Ambos contenedores se ejecutar simulataneamente.
% - Este experimento se ejecutar un total de 10 veces y de los resultados obtenidos se scara la media respecto al numero total de ejecucion para cada herramienta. Cabe
% destacar que se han dejado 2 minutos entre un tanda de experimentos y otra para evitar que el calor de la CPU influya en el consumo energetico influya en los resultados
% finales de las pruebas. 
% - Una vez finalizadas las pruebas y llevado a cabo el post-procesamiento indicado con anterioridad, se observa que ninguna de lso dos herramitnas es capaz de medir
% correctamente la energia del nodo y la energia de los dos contendores, denotando una clara falta de energia respecto al consumo total del sistema, el cual se ha medido
% empleadno la herramienta EM. 


% +++ Uso de Herramiena Propia +++
% ++++++ EXPERIMENTO 4 +++++++++
% - Debido a los insatisfactorios resultados obtenidos con las herramientas kepler y scp, se ha propuesto utilizar una herramientas propia desarrollada en el proyecto Eco-Stream,
%   la cual hace uso de diferentes herramientas y aplicaciones internas de los sistemas linux para lograr una precision mucho mayor a la hora de obtener metricas de consumo.
% - Para comprobar su efectividad, se llevaran a cabo los mismos experimentos aplicando las mismas limitaciones sobre los conteneodres.
% - Comparando los datos obtenidos, se puede observar que la herramienta MONITORING es mucho mas efectiva, tanto en la estimacion de consumo del nodo como en la toma de metricas multiples procesos ejecutandose simultanemente. Los resultados obtenido de la experimentacion se pueden observar a continuacion en la tabla REF TABLA
% 
\subsection{Conclusiones fase 1 de la experimentacion }
% ++++++ Python Plot & POST PROCESADO DATOS ++++++++++++
% + Comentar como se han procesado los datos y de devolvia cada herrameinta y enque forma
%   -Formato datos Devueltos:
%       - Kepler --> Informacion resumida desde el inicio hasta la peticion de datos
%       - Schapndre --> Metricas tomadas cada X segundos(En nuestro caso se define 1), cada bucle metrica ssitema en ese instante
%       - EM ---> JSON con infromacionconsumo (Tiempo, pkg, d_ram, total)
%       - MONITOR :
%           - Informacion resumida en .global
%           - Informacion resumida en .csv cada iteracion(Igual que scaphandre)
%           - Consumo total del sistema basandose en EM 
%   - Procesado de los datos 
%       - Bash script para schaphandre 
%           - Caracteristicas XXX 
%       - Bash script kepler 
%           - Caracteristicas X 
%       - Bash para METRICS 
%           - Caracteriticas X
% - En las siguientes graficas se aprecia la medida de los datos obtenidos con cada herramienta en cada experiménto, y como se puede aprecias la herramienta MONITORING es la 
%   mas adecuada para la fase 2 del proyecto.
% - Se ha utilizado python con seaborn y matplotlib para crear las tablas comparativas de datos. Ha sido necesario instalar las 
%  librerias pyqt6(QUE ES ?? ) y la libreria de linux que se encarga de poder plotear figuras flotantes libxcb-cursor0